\documentclass[10pt]{article}
\usepackage[letterpaper,margin=0.62in]{geometry}
\usepackage{booktabs,array,amsmath,amssymb,xcolor,enumitem}
\usepackage[hidelinks]{hyperref}
\setlength{\parindent}{0pt}
\setlength{\parskip}{4pt}
\setlist{nosep,leftmargin=*}
\newcommand{\cafa}{\textbf{CAFA-IVR}}
\begin{document}
\begin{center}
{\LARGE\bfseries CAFA-IVR v1.0}\\[3pt]
{\large Counterfactual ASR Failure Attribution for Conversational IVR}\\[6pt]
\textbf{Reference Specification} \quad | \quad Vijay Kumar Sridharan \quad | \quad August 2026\\[2pt]
IEEE Senior Member \quad | \quad Frisco, TX, USA
\end{center}
\hrule

\section*{1. Purpose}
\cafa{} defines a portable test procedure for determining whether adding an audio/ASR layer causes a conversational test to fail when the same semantic request succeeds through a reference-text control. The framework is independent of ASR vendor, NLU vendor, contact-center platform, and dialogue architecture.

\section*{2. Required paired execution}
For every semantic test $i$, an implementation \textbf{MUST} execute both paths using the same frozen downstream configuration:
\[
T_i:\ \text{reference text}\rightarrow\text{downstream system},\qquad
A_i:\ \text{audio}\rightarrow\text{ASR}\rightarrow\text{downstream system}.
\]
The two paths \textbf{MUST} share the same expected semantic outcome. The downstream component may be an intent classifier, workflow, LLM agent, or tool-using system, provided a reproducible task oracle exists.

\section*{3. Attribution states}
\begin{center}
\begin{tabular}{cclp{2.7in}}
\toprule
Text & Audio & State & Interpretation\\
\midrule
Pass & Pass & \texttt{HEALTHY} & End-to-end semantic success.\\
Pass & Fail & \texttt{SPEECH\_ATTRIBUTABLE} & Failure appears only after the speech path is introduced.\\
Fail & Fail & \texttt{DOWNSTREAM\_OR\_TEST} & Reference text already fails; do not attribute the failure to ASR.\\
Fail & Pass & \texttt{CONTEXT\_OR\_ORACLE} & Investigate nondeterminism, context, or the expected-result oracle.\\
\bottomrule
\end{tabular}
\end{center}
A speech-attributable state narrows the failure surface; it does not by itself prove one acoustic root cause.

\section*{4. Metrics}
\textbf{WER.} Preserve conventional lexical fidelity:
\[
\mathrm{WER}=\frac{S+D+I}{N}.
\]

\textbf{ASR-IFR.} Let $F_i=1$ when the text control passes and the audio-path task fails, else $0$:
\[
\mathrm{ASR\mbox{-}IFR}=\frac{1}{N}\sum_i F_i.
\]
A conditional variant may divide by the number of text-control successes.

\textbf{CEER.} Critical-Entity Error Rate measures governed semantic values such as amounts, dates, negation, account type, authentication values, and lock/unlock/cancel/confirm actions. Entity values \textbf{SHOULD} be normalized before scoring.

\textbf{CIER (optional).} A deployment may apply severity weights to L0--L4 outcomes. The reference scorer uses illustrative weights L0=0, L1=0, L2=.25, L3=.65, L4=1.0. These are policy examples, not universal thresholds.

\section*{5. Minimum trial record}
\begin{center}\small
\begin{tabular}{lp{1.0in}p{4.4in}}
\toprule
Field & Status & Purpose\\
\midrule
\texttt{test\_id} & Required & Stable semantic test identifier.\\
\texttt{condition} & Required & Clean, telephony, noise, rate, or another controlled condition.\\
\texttt{reference\_text} & Required & Canonical human-authored reference.\\
\texttt{asr\_transcript} & Required & Observed speech-engine transcript.\\
Expected/observed outcome & Required & Intent labels or equivalent task oracle.\\
\texttt{text\_control\_pass} & Required & Downstream result when ASR is bypassed.\\
\texttt{audio\_task\_success} & Required & Outcome through audio $\rightarrow$ ASR.\\
Entity JSON & Recommended & Governed critical values and observed values.\\
\texttt{risk\_tier} & Recommended & Local Low/Medium/High or equivalent.\\
Versions & Recommended & ASR, NLU/agent, prompt/tool/oracle configuration.\\
\bottomrule
\end{tabular}
\end{center}

\section*{6. Minimum acoustic matrix}
A production suite \textbf{SHOULD} establish clean/wideband baseline, 8-kHz or actual telephony path, moderate and severe controlled noise, bounded speaking-rate variation, representative human speakers, and critical-semantic cases. Synthetic speech may expand coverage but should not be the only evidence used for a human-robustness claim.

\section*{7. Release-gating policy}
\cafa{} does not prescribe universal numeric thresholds. A team should compare a candidate against its last approved baseline and use stricter policy for high-consequence journeys. A practical policy is:
\begin{enumerate}
\item no unexplained meaningful increase in ASR-IFR;
\item no new governed critical-entity failures in high-risk journeys without explicit approval;
\item WER regressions are investigated, but WER alone does not decide the release;
\item text-control failures are routed to NLU/agent/test owners rather than ASR owners;
\item failed gates retain audio, transcript, versions, expected outcome, and observed outcome for replay.
\end{enumerate}

\section*{8. CI/CD lifecycle}
\begin{center}
\fbox{Freeze versions} $\rightarrow$ \fbox{Run text controls} $\rightarrow$ \fbox{Run audio suite} $\rightarrow$ \fbox{Score}\\[4pt]
$\rightarrow$ \fbox{Compare baseline} $\rightarrow$ \fbox{Pass / attributed investigation / rerun}
\end{center}

An adapter only needs to populate the trial record. Credentials should stay in the adopting platform's standard credential mechanism and must not be committed to a CAFA results package.

\section*{9. Reproducibility and independent adoption}
An auditable run should freeze the audio or source identifiers, perturbation parameters and random seeds, ASR version, downstream version, prompts/tools when applicable, locale, endpoint configuration, test-oracle version, and CAFA-IVR specification version.

An organization may accurately state that it adopted CAFA-IVR only when it actually used the paired procedure. A useful implementation statement is:

\begin{quote}\small
``We evaluated conversational voice regressions using CAFA-IVR v1.0 paired reference-text and audio-path testing and incorporated ASR-IFR into our failure-triage workflow.''
\end{quote}

Independent adopters should retain test scope, decision, resulting engineering change, and artifacts sufficient for a third party to verify the use.

\section*{10. Conformance checklist}
A CAFA-IVR v1.0 conforming implementation should be able to answer \emph{yes} to the following:
\begin{itemize}
\item Each semantic test has a stable ID, reference text, and expected outcome.
\item Reference text and audio use the same frozen downstream configuration.
\item The ASR transcript and both path outcomes are retained at trial level.
\item ASR-IFR is calculated from the paired outcomes, not inferred from WER.
\item High-risk critical entities are explicitly annotated or the absence of CEER is documented.
\item Candidate results are compared with a named approved baseline when used for release gating.
\item Failed trials retain enough evidence for replay and independent review.
\item The CAFA-IVR specification version and tested ASR/downstream versions are recorded.
\end{itemize}

\section*{11. Independent adoption record}
When the framework is evaluated or adopted by an independent team, retain the following fields so use can be verified later:
\begin{center}\small
\begin{tabular}{lp{4.9in}}
\toprule
Field & Record\\
\midrule
Organization / team & Name and business unit\\
Evaluator & Name, title, and contact\\
CAFA-IVR version & Specification/release used\\
Technical scope & ASR, NLU/agent, number of tests, number of audio trials\\
Evidence type & Human, synthetic, or mixed speech; acoustic conditions\\
Operational use & Test plan, CI job, dashboard, runbook, release decision, or evaluation only\\
Engineering impact & Configuration/model change, avoided misdiagnosis, new gate, issue found, release decision\\
Evidence retained & Report, screenshot, commit, run ID, URL, artifact hash, or internal document reference\\
Public citation permission & Yes / No / Limited\\
\bottomrule
\end{tabular}
\end{center}

\section*{12. Non-endorsement}
CAFA-IVR is a technical reference framework. Listing an employer, vendor, dataset, or platform does not imply sponsorship or endorsement. Production thresholds and risk decisions remain the responsibility of the adopting organization.

\end{document}
